% !TeX root = ../thuthesis-example.tex

\begin{comments}
% \begin{comments}[name = {指导小组评语}]
% \begin{comments}[name = {Comments from Thesis Supervisor}]
% \begin{comments}[name = {Comments from Thesis Supervision Committee}]

  针对中性粒细胞在免疫效应执行过程中具有明显情境依赖性、其关键功能依赖蛋白尚缺乏系统识别这一问题，该论文围绕免疫情境中的候选必需蛋白筛选与验证开展了较为系统和深入的研究。《基于深度学习的中性粒细胞必需蛋白筛选与验证研究》首先基于蛋白语言模型构建了Human-level与Immune-level双层蛋白必需性评估框架，采用ProtT5提取序列特征，并结合注意力机制建立了必需性预测模型。其次，论文基于protein essential score（PES）及其差异指标筛选免疫情境候选必需蛋白，并结合理化特征、功能模块、注意力图谱和公共数据库证据，对候选蛋白的生物学合理性进行了较为全面的分析。最后，论文围绕ATP6V1B2等代表性候选蛋白开展了机制解析及体外功能验证，结果表明ATP6V1B2在氧化爆发-颗粒/溶酶体高负荷活化状态下具有更突出的情境增强依赖特征。

  该论文选题紧密结合免疫调控与靶点发现中的前沿问题，具有较好的理论意义和一定的应用价值。论文在问题凝练、方法设计和结果论证方面均体现出较强的系统性，提出的研究思路具有一定创新性，为发现具有细胞选择性的免疫调控靶点提供了新的技术路径。论文作者能够较好地综合运用生物医学工程、深度学习和实验研究等方面的知识开展工作，表现出较好的科研素养、独立开展研究的能力以及分析和解决问题的能力。论文技术路线清晰，研究设计合理，实验结果和分析结论具有较好的可信度。

  论文结构完整，层次清楚，逻辑严谨，文字表达规范，符合硕士学位论文写作要求。综上所述，该论文已达到工学硕士学位论文水平，同意唐健同学参加硕士学位论文答辩。

\end{comments}
